\documentclass[11pt,compress,t,notes=noshow, xcolor=table]{beamer}

\input{../../style/preamble}
\input{../../latex-math/basic-math}
\input{../../latex-math/basic-ml}
\input{../../latex-math/optim-basics}

\title{Optimization in Machine Learning}

\begin{document}

\titlemeta{
Constrained Optimization
}{
Primal Coordinate Descent
}{
figure/primal_cd_path.pdf
}{
\item Coordinate-wise minimization over feasible sets
\item Convex slices reduce to interval-constrained subproblems
\item Box regression and truncation / clipping updates
\item When constrained coordinate descent is attractive
}

\begin{framev}[fs=small]{Why Coordinate Descent Fits Constraints}
\splitVTT[0.46]{
\textbf{Problem}
$$
\min_{\wv \in \S} F(\wv).
$$

\begin{itemize}
\item A full gradient step can violate several constraints at once.
\item Fixing the other coordinates reduces feasibility to a one-dimensional slice of $\S$.
\item For convex $\S$, that slice is an interval and therefore easy to search.
\end{itemize}
}{
\imageC[0.65]{figure/primal_cd_intervals.pdf}
}
\end{framev}


\begin{framei}[fs=small]{Constrained Coordinate Descent}
\item Start from a feasible point $\wv^{[0]} \in \S$ and choose the active coordinate $i_t$ cyclically or at random.
\item For the current iterate, compute the feasible slice
$$
I_t := \left\{ z \in \R \mid (\wv_{-i_t}^{[t]}, z) \in \S \right\},
$$
where $\wv_{-i_t}^{[t]}$ denotes all coordinates except the active one.
\item Solve the one-dimensional constrained problem
$$
w_{i_t}^{[t+1]} \in \argmin_{z \in I_t} F(\wv_{-i_t}^{[t]}, z).
$$
\item Keep all other coordinates unchanged:
$$
w_j^{[t+1]} = w_j^{[t]} \qquad \forall j \neq i_t.
$$
\item One full pass over all coordinates is a \textbf{sweep}; \textbf{block coordinate descent} updates groups instead of single variables.
\end{framei}


\begin{framei}[fs=small]{Convex Sets Give Interval Subproblems}
\item For a feasible point $\wv$ and coordinate $i$, define the slice
$$
I_i(\wv) := \left\{ z \in \R \mid (w_1,\dots,w_{i-1},z,w_{i+1},\dots,w_d) \in \S \right\}.
$$
\item Geometrically, $I_i(\wv)$ is the intersection of $\S$ with an affine line parallel to axis $i$.
\item If $\S$ is convex, then $I_i(\wv)$ is a convex subset of $\R$, hence an interval.
\item Therefore, if $F$ is convex, each coordinate step solves a \textbf{univariate convex problem over an interval}.
\item The hard part is no longer a full $d$-dimensional update, but computing the interval endpoints and minimizing along that slice.
\item This is why coordinate descent is especially appealing when the constraints have simple coordinate-wise structure.
\end{framei}


\begin{framei}[fs=small]{Worked Example: One Coordinate Step}
\item Consider optimizing $F(w_1,w_2,w_3)$ subject to
$$
w_1^2 - w_1 w_2 + \frac{w_2^2}{4} + 3 w_2 w_3 + 4 w_3^2 \le 4,
$$
$$
2 w_1 + w_2 - 3 w_3 \le 4.
$$
\item In a coordinate-descent step, fix $w_2 = 2$ and $w_3 = 0$, and optimize only $w_1$.
\item Plugging these values into the constraints gives
$$
w_1^2 - 2 w_1 - 3 = (w_1 - 3)(w_1 + 1) \le 0,
\qquad
w_1 \le 1.
$$
\item Hence
$$
w_1 \in [-1,3] \cap (-\infty,1] = [-1,1].
$$
\item The subproblem becomes
$$
\min_{w_1 \in [-1,1]} G(w_1),
\qquad
G(w_1) := F(w_1,2,0).
$$
\end{framei}


\begin{framei}[fs=small]{How Do We Solve the 1D Subproblem?}
\item Let $G(z)$ denote the one-dimensional objective over a feasible interval $[a,b]$.
\item If $G$ is differentiable and convex, first solve the unconstrained stationarity equation
$$
G'(z) = 0.
$$
\item If the solution $\hat z$ lies in $[a,b]$, it is the coordinate update.
\item Otherwise, the optimum lies at one of the endpoints, so the exact update is
$$
z^\ast = \Pi_{[a,b]}(\hat z).
$$
\item If no closed form exists, one can still solve the univariate problem cheaply by line search or a few Newton steps.
\item Coordinate descent therefore replaces one hard multivariate problem by many cheap one-dimensional problems.
\end{framei}


\begin{framev}[fs=small]{Box Constraints Are Especially Friendly}
\splitVTT[0.50]{
\textbf{Feasible set}
$$
\S = \left\{\wv \in \R^d \mid
\begin{array}{l}
l_j \le w_j \le u_j \\
j = 1,\dots,d
\end{array}
\right\}.
$$

\textbf{Coordinate slice}
$$
I_i(\wv) = [l_i, u_i].
$$

\begin{itemize}
\item The interval does not depend on the other coordinates.
\item In 2D, coordinate descent moves along horizontal and vertical segments.
\item Every step stays feasible without a global projection.
\end{itemize}
}{
\imageC[0.92]{figure/primal_cd_path.pdf}
}
\end{framev}


\begin{framei}[fs=small]{ML Application: Box Regression}
\item Add lower and upper bounds to ridge-style linear regression:
$$
\min_{\wv} J(\wv) = \frac{1}{2}\|\mathbf{D}\wv - \yv\|_2^2 + \frac{\lambda}{2}\|\wv\|_2^2
$$
$$
\text{subject to } \qquad l_i \le w_i \le u_i, \qquad i=1,\dots,d.
$$
\item $\mathbf{D} \in \R^{n \times d}$ is the feature matrix, $\wv \in \R^d$ the coefficient vector, and $\yv \in \R^n$ the response vector.
\item For fixed coordinates except $w_i$, the objective is a simple univariate quadratic over the interval $[l_i,u_i]$.
\item Hence box regression can be solved by coordinate descent with exact closed-form coordinate updates.
\item \textbf{Special case}: nonnegative least squares uses $l_i = 0$ and no finite upper bound.
\end{framei}


\begin{framev}[fs=small]{Truncated Coordinate Update for Box Regression}
\splitVTT[0.58]{
Let $\mathbf{d}_i$ be column $i$ of $\mathbf{D}$ and let
$$
\mathbf{r} = \yv - \mathbf{D}\wv
$$
be the current residual vector.

\textbf{Unconstrained update}
$$
w_i \leftarrow
\frac{w_i \|\mathbf{d}_i\|_2^2 + \mathbf{d}_i^\top \mathbf{r}}
{\|\mathbf{d}_i\|_2^2 + \lambda}
$$

\textbf{Constrained update}
$$
w_i \leftarrow
T_i\!\left(
\frac{w_i \|\mathbf{d}_i\|_2^2 + \mathbf{d}_i^\top \mathbf{r}}
{\|\mathbf{d}_i\|_2^2 + \lambda}
\right)
$$
$$
T_i(x) = \min\{u_i, \max\{l_i, x\}\}
$$
}{
\imageC[0.90]{figure/primal_cd_truncation.pdf}
}
\end{framev}


\begin{framei}[fs=small]{Practical Remarks}
\item Coordinate descent is attractive when each slice $I_i(\wv)$ is easy to compute and the one-dimensional minimizer is cheap.
\item Box, nonnegativity, and other separable constraints are ideal because each coordinate update only needs a clipping step.
\item For convex smooth objectives, cyclic or randomized coordinate descent converges; strong convexity often yields linear rates.
\item Block coordinate descent can be preferable when variables are strongly coupled.
\item Limitations remain important: for general coupled constraints, computing the slice itself may already be expensive.
\item Coordinate descent can also stall at coordinate-wise minima in more difficult nonsmooth or nonconvex settings.
\end{framei}


\begin{framev}{Take-Away}
\begin{enumerate}
\item Fix all coordinates except one.
\item Intersect the feasible set with that coordinate line.
\item Solve the resulting one-dimensional problem over an interval.
\item Sweep cyclically, randomly, or in blocks.
\end{enumerate}

\vfill
\oneliner{Primal coordinate descent is most useful when constraint structure makes each coordinate slice simple, especially for box-type constraints.}

\vfill
Next, we switch from primal feasible updates to Lagrangian relaxation and duality.
\end{framev}

\endlecture
\end{document}
